\section*{Chapter \ref{ch:b2:lab-techniques-2} --- Lab Techniques II: Synthesis and Analysis in Practice}

\begin{solution}{exo:b2:lab-techniques-2:1}
Ice and water; ice and salt; dry ice and acetone (as for the \omterm{def:b2:enolates-aldol:enolate}{enolates} of \cref{ch:b2:enolates-aldol}).
\end{solution}

\begin{solution}{exo:b2:lab-techniques-2:2}
The lower layer. Add a few drops of water: they join the aqueous layer, here the upper one.
\end{solution}

\begin{solution}{exo:b2:lab-techniques-2:3}
The first portions clump as they take up water; when newly added solid stays loose and swirls freely,
the solution is dry.
\end{solution}

\begin{solution}{exo:b2:lab-techniques-2:4}
$R_f$ is too high for a column (about 0.3 is wanted) and the two spots are close: use a less polar
eluent, more hexane (9\,:\,1, for instance), which lowers both $R_f$ and widens their separation in
column volumes.
\end{solution}

\begin{solution}{exo:b2:lab-techniques-2:5}
\begin{itemize}
  \item Sodium hydrogencarbonate (pH 8.3 $>$ 4.2 + 2): benzoate in water; acidify, extract benzoic acid.
  \item Sodium hydroxide (pH $>$ 12 $>$ 9.99 + 2): phenoxide in water; acidify, extract phenol.
  \item Hydrochloric acid (pH about 1 $<$ 4.6 $-$ 2): anilinium in water; make basic, extract aniline.
  \item Naphthalene stays in the ether: dry, evaporate.
\end{itemize}
The hydrogencarbonate wash must come first: hydroxide would extract the acid and the phenol together.
% ledger: ph:NaHCO3, pka:benzoic, pka:phenol, pka:anilinium
\end{solution}

\begin{solution}{exo:b2:lab-techniques-2:6}
$1/T = 1/T_b - (R/\Delta_{\mathrm{vap}}H)\ln(p/p^\circ)$: bromobenzene about \qty{49}{\degreeCelsius},
benzaldehyde about \qty{68}{\degreeCelsius} (figure of \cref{prop:b2:lab-techniques-2:reduced-pressure}).
\end{solution}

\begin{solution}{exo:b2:lab-techniques-2:7}
$0.050 \times 200\,000/400 = \qty{25}{K}$.
\end{solution}

\begin{solution}{exo:b2:lab-techniques-2:8}
Each compound has its own response: area per cent is a mass per cent only if the \omterm{def:b2:chromatography:quantitative}{response factors} are
equal. NMR: $0.12/3 = 0.040$~mol of impurity per mol of product, that is $0.040 \times 100 = \qty{4}{g}$
per \qty{200}{g} of product; $w = 200/204 \approx 98.0~\%$.
\end{solution}

\begin{solution}{exo:b2:lab-techniques-2:9}
Theoretical mass $0.0300 \times 164 = \qty{4.92}{g}$. Uncorrected $4.10/4.92 \approx 83.3~\%$; corrected
$0.950 \times 4.10/4.92 \approx 79.2~\%$.
\end{solution}

\begin{solution}{exo:b2:lab-techniques-2:10}
Benzene means the reagent formed and was then protonated: water in the glassware, the solvent or the
aldehyde; air (moisture) entering through a leak; an aldehyde partly oxidised to benzoic acid, whose
acidic hydrogen destroys the reagent. Remedies: oven-dried glassware, dry solvent, freshly distilled
aldehyde, a nitrogen overpressure checked at the bubbler. Little reagent at all would point to an
initiation failure (a magnesium surface covered by oxide), avoided by fresh or activated turnings.
\end{solution}

\begin{solution}{exo:b2:lab-techniques-2:11}
\qty{270}{kJ} at \qty{150}{W}: at least $270\,000/150 = \qty{1800}{s}$, half an hour, and only if the
halide reacts as fast as it is added. Faster, the heat release exceeds the cooling: the temperature rises,
or reagent accumulates and may later react all at once (\cref{prop:b2:lab-techniques-2:accumulation}).
\end{solution}

\begin{solution}{exo:b2:lab-techniques-2:12}
Recrystallisation: $13.0 \times 0.80 = \qty{10.4}{g}$ at 99~\%. Column: $13.0 \times 0.95 =
\qty{12.4}{g}$ at 99.5~\%, but \qty{1.5}{L} of solvent to buy, evaporate and dispose of. The column gives
about \qty{2}{g} more of slightly purer product; recrystallisation is simpler, cheaper and less
wasteful, and is preferred when its purity suffices.
\end{solution}

\begin{solution}{pb:b2:lab-techniques-2:1}
\textbf{1.} Ethoxyethane: extremely flammable, harmful, forms peroxides. Bromobenzene: flammable, toxic,
toxic to aquatic life. Magnesium: flammable solid, releases hydrogen with water.
\textbf{2.} \ce{C6H5MgBr + H2O -> C6H6 + Mg(OH)Br}: every drop of water destroys reagent; oven drying
removes the film of water on the glass.
\textbf{3.} To keep out moisture and oxygen, which also destroys the reagent.
\textbf{4.} The flow of nitrogen; it prevents air from entering through the outlet.
\textbf{5.} It does not react with the reagent, its oxygen lone pairs stabilise the magnesium, and its
low boiling point (\qty{307.7}{K}) lets reflux limit the temperature.
\textbf{6.} $15.7/157.01 = \qty{0.1000}{mol}$ bromobenzene; $2.67/24.305 = \qty{0.1099}{mol}$ magnesium,
in 10~\% excess.
\textbf{7.} $0.1000 \times 270 = \qty{27.0}{kJ}$.
\textbf{8.} So that the heat is released at the rate at which reflux and the bath remove it.
\textbf{9.} $0.0200 \times 270 = \qty{5.40}{kJ}$.
\textbf{10.} $5400/170 \approx \qty{32}{K}$.
\textbf{11.} About \qty{52}{\degreeCelsius}, far above the boiling point of ethoxyethane
(\qty{34.6}{\degreeCelsius}): violent boiling that can overwhelm the condenser and throw flammable vapour
out of the flask.
\textbf{12.} Add about a tenth of the bromobenzene, wait for the reaction to start (cloudiness, gentle
reflux), then add the rest at the rate of a gentle reflux, with the ice bath ready.
\textbf{13.} \ce{C6H5MgBr + C6H5CHO} gives the alkoxide \ce{(C6H5)2CHOMgBr}; benzaldehyde,
$9.55/106.12 = \qty{0.0900}{mol}$, is limiting.
\textbf{14.} The alcohol is benzylic twice over: a strong acid would turn it into a stabilised cation,
leading to ethers or substitution products; ammonium chloride is acidic enough to protonate the alkoxide
and dissolve the magnesium salts.
\textbf{15.} The upper layer: ethoxyethane is less dense than water.
\textbf{16.} It removes most of the water dissolved in the ether.
\textbf{17.} Magnesium sulfate until free-flowing, filtration, then a rotary evaporator under reduced
pressure.
\textbf{18.} Phenylmagnesium bromide couples with unreacted bromobenzene; a high local concentration of
bromobenzene and high temperature favour it, another reason for a slow addition.
\textbf{19.} Biphenyl ($R_f$ 0.85), the less polar.
\textbf{20.} About $1/0.85 \approx 1.2$ column volumes for biphenyl, $1/0.25 = 4$ for the alcohol.
\textbf{21.} The alcohol contributes 10 \omterm{def:b2:huckel:aromatic}{aromatic} protons per CH; the extra $10.90 - 10.00 = 0.90$ comes from
biphenyl, 10 protons per molecule: 0.090 mol of biphenyl per mol of alcohol.
\textbf{22.} $w = 184.24/(184.24 + 0.090 \times 154.21) \approx 0.930$.
\textbf{23.} $12.50 \times 0.930 \approx \qty{11.62}{g}$.
\textbf{24.} $0.0900 \times 184.24 \approx \qty{16.58}{g}$.
\textbf{25.} $12.50/16.58 \approx 75.4~\%$.
\textbf{26.} $11.62/16.58 \approx \boldsymbol{70~\%}$ (70.1~\%), purity-corrected.
\end{solution}
